\section[Intersection relations]{Actions on posets and their intersection relations}
\label{sec-osa-intersection}

An action is \emph{ordered by elements} if, for parallel interspaces
$f,g:X\to Y$, the inequalities $\supfun{f}x\sqsubseteq\supfun{g}x$
for every $x\in P_X$ imply $f\sqsubseteq g$. In particular the action
then distinguishes parallel interspaces. In a \emph{functional} action
the interspaces themselves are the monotone maps that act on the posets.
Every action induces a functional action by taking, for each $X,Y$,
the image of $\Hom(X,Y)$ under $f\mapsto\supfun{f}$ with pointwise
order. Keep the objects $X$ as labels for the posets $P_X$, even when
two of the posets coincide. Composition of image maps is well defined
because the original action preserves composition. An action ordered
by elements is isomorphic to this functional action.

An \emph{ordered semifunctor} is a semifunctor that is monotone on each
Hom-set. If $F:\mathcal D\to\mathcal C$ is an ordered semifunctor and
$\mathcal C$ acts on the posets $P_X$, the \emph{pullback action} along
$F$ assigns to an object $U$ of $\mathcal D$ the poset $P_{F(U)}$ and
to an arrow $f$ the map $\supfun{F(f)}$. The two composition laws and
the two monotonicity conditions follow from those for $F$ and the given
action.

I use the ordered dagger semicategories of
Section~\osaprereqref{sec-ordered-dagger}. In particular, a dagger
takes $f:X\to Y$ to $f^\dagger:Y\to X$, preserves the order and
reverses composition. An \emph{ordered dagger semicategory action}
is an ordered semicategory action whose underlying semicategory has
such a dagger. In the one-object case this is an ordered semigroup
action with an involution on the semigroup. Compatibility of the two
action maps associated to $f$ and $f^\dagger$ can be expressed by
intersection.

\begin{defn}\label{def-osa-relation}
\index{intersection relation}
The \emph{intersection relation} $\mathord{\suprel{f}}\subseteq P_X\times P_Y$
of an interspace $f:X\to Y$ is
\[
 x\suprel{f}y\quad\Longleftrightarrow\quad y\nasymp\supfun{f}x
 \qquad(x\in P_X,\ y\in P_Y).
\]
An ordered dagger semicategory action is \emph{intersection-symmetric}
if $\mathord{\suprel{f^\dagger}}=\mathord{\suprel{f}}^{-1}$ for every $f$,
where the inverse of a relation is its converse. Equivalently,
$y\nasymp\supfun{f}x$ iff $x\nasymp\supfun{f^\dagger}y$.
\end{defn}

\begin{prop}\label{prop-osa-relation}
For every ordered action, $f\sqsubseteq g$ implies
$\mathord{\suprel{f}}\subseteq\mathord{\suprel{g}}$. If $P_Y$ is separable
in the sense of Section~\osaprereqref{sep-and-full}, $\suprel{f}$
determines $\supfun{f}$. If $P_Y$ is strongly separable and the action is
ordered by elements on $\Hom(X,Y)$, then
$\mathord{\suprel{f}}\subseteq\mathord{\suprel{g}}$ implies $f\sqsubseteq g$.
\end{prop}

\begin{proof}
The set of $y$ related to $x$ is the full star of $\supfun{f}x$.
The full star is monotone. Separability makes it injective, and
strong separability makes it order reflecting. Apply these facts for
each $x$ and use order by elements for the last assertion.
\end{proof}

\begin{prop}\label{prop-osa-rel-comp}
For $f:X\to Y$ and $g:Y\to Z$,
$\mathord{\suprel{g\circ f}}=\mathord{\suprel{g}}\circ\supfun{f}$,
where a function in a relational composition stands for its graph.
In an intersection-symmetric action, for $x\in P_X$ and $z\in P_Z$,
\[
 x\suprel{g\circ f}z
 \quad\Longleftrightarrow\quad
 \supfun{f}x\nasymp\supfun{g^\dagger}z.
\]
If $P_Y$ is atomic with least element (Section~\osaprereqref{sec-common-atoms}),
this is equivalent to the
existence of an atom $y$ of $P_Y$ with $x\suprel{f}y$ and $y\suprel{g}z$.
\end{prop}

\begin{proof}
The composition law for the action gives
$z\nasymp\supfun{g\circ f}x$ iff $\supfun{f}x\suprel{g}z$.
Apply intersection symmetry for the second formula. In an atomic
poset with least element, two elements intersect iff an atom lies
below both. An atom intersects an element iff it lies below that
element, giving the last assertion.
\end{proof}

The following construction gives actions whose arrows record the forward
and reverse maps explicitly. For posets $P,Q$ take pairs of maps
$\alpha:P\to Q$ and $\beta:Q\to P$ satisfying
\begin{equation}\label{eq-osa-compatible-pair}
 y\nasymp\alpha x\quad\Longleftrightarrow\quad x\nasymp\beta y
 \qquad(x\in P,\ y\in Q).
\end{equation}
Composition applies the forward maps in order and the reverse maps in
reverse order. Indeed, if $(\gamma,\delta)$ is a second pair from $Q$ to
a poset $R$, then for $z\in R$,
$z\nasymp\gamma\alpha x$ iff $\alpha x\nasymp\delta z$ iff
$x\nasymp\beta\delta z$. Identity maps form the identity pair and
interchanging the two maps gives the dagger. These operations give a
dagger category with the posets as objects. If the objects are strongly
separable, both maps of every pair are monotone. For example,
$x\sqsubseteq x'$ implies that every $y$ intersecting $\alpha x$ also
intersects $\alpha x'$, by \eqref{eq-osa-compatible-pair}. Hence
$\alpha x\sqsubseteq\alpha x'$. Interchanging the maps proves monotonicity
of $\beta$. With the componentwise pointwise order on pairs, composition
is monotone and evaluation of the forward map is an intersection-symmetric
ordered action. Strong separability also gives order by elements. An
inequality between forward maps gives inclusion of the corresponding
relations and hence inclusion of the full stars of the reverse values
at each $y$, which gives the inequality between reverse maps.

If $P,Q$ are bounded, the least compatible pair consists of the
two constant least maps. The greatest compatible pair sends each
least input to the least output and each non-least input to the
greatest output. For the latter pair both sides of
\eqref{eq-osa-compatible-pair} hold exactly when $x$ and $y$ are
both non-least. These are the bounds in the componentwise order:
any compatible pair preserves the least elements, since the image
of a least element intersects no element by
\eqref{eq-osa-compatible-pair}.

\begin{prop}\label{prop-osa-core-recovery}
For each of two action posets $P_X,P_Y$ choose a subset $Z_X,Z_Y$,
respectively. For an element $x\in P_X$, put
$\up_{Z_X}x=\setcond{u\in Z_X}{x\sqsubseteq u}$, and use the same
notation for elements of $P_Y$. Suppose that in each poset an
element intersects a fixed element iff all its upper approximants
in the selected subset intersect that fixed element. In an
intersection-symmetric action one then has
\[
 x\suprel{f}y
 \iff
 \forall(u,v)\in \up_{Z_X}x\times\up_{Z_Y}y:\ u\suprel{f}v.
\]
This formula applies also to a pair satisfying
\eqref{eq-osa-compatible-pair}. With the hypothesis on $Z_X$ alone,
one may approximate $x$ alone. If $P_Y$ is separable and
intersection with $y\in P_Y$ commutes with the meet of the
values $\supfun{f}u$ for $u\in\up_{Z_X}x$, then
\[
 \supfun{f}x=\bigsqcap_{u\in\up_{Z_X}x}\supfun{f}u.
\]
Here commuting means that a test element intersects the meet iff
it intersects every member of the family, and requires that the
meet exist.
\end{prop}

\begin{proof}
Use the compatibility equation to test $x$ against
$\supfun{f^\dagger}y$. The approximation hypothesis for $x$ gives
the test at each $u\in\up_{Z_X}x$. For each $u$, approximate $y$
in the test of its intersection with $\supfun{f}u$. If the meet
satisfies the stated condition, it has the same full star as
$\supfun{f}x$. Separability gives equality.
\end{proof}

\begin{prop}\label{prop-osa-atom-recovery}
Suppose $P_X$ and $P_Y$ are atomic with least elements.
In an intersection-symmetric action $x\suprel{f}y$ holds iff
there are atoms $p\sqsubseteq x$ and $q\sqsubseteq y$ with
$p\suprel{f}q$.
If both posets are atomistic and the action is ordered by elements,
the relation on atoms determines $f$ and reflects its order.
For every $x\in P_X$ one has
$\supfun{f}x=\bigsqcup_{p\in\atoms^{P_X},\,p\sqsubseteq x}\supfun{f}p$.
\end{prop}

\begin{proof}
An atom $q$ below $y$ and $\supfun{f}x$ exists exactly when these
elements intersect. Equivalently, $x$ intersects $\supfun{f^\dagger}q$,
so choose an atom $p$ below both. The converse uses monotonicity.
If every element is a join of its atoms, the action at $x$ is
the join of the atoms $q$ obtained by this test. This also shows
that it is the least upper bound of the values at the atoms below $x$.
Inclusion of the relation on atoms thus gives pointwise order of
the action values and hence arrow order.
\end{proof}

\section{Images, domains, and preservation of joins}
\label{sec-osa-image}

\begin{defn}\label{def-osa-image}
\index{image!of an interspace}\index{domain!of an interspace}
For $f:X\to Y$ with $P_X$ having a greatest element $\top^{P_X}$,
define the \emph{image} $\im f=\supfun{f}\top^{P_X}\in P_Y$.
In a dagger action with $P_Y$ having a greatest element, define the
\emph{domain} $\dom f=\im(f^\dagger)\in P_X$.
\end{defn}

Both are monotone on their Hom-sets. Every $\supfun{f}x$ lies below
$\im f$ when it is defined. For $g:Y\to Z$ one has
$\im(g\circ f)=\supfun{g}(\im f)$, by the composition law.

\begin{prop}\label{prop-osa-support}
In an intersection-symmetric action, let $f:X\to Y$ with $P_Y$ having
a greatest element. For each $x\in P_X$, $x\nasymp\dom f$ iff
$\supfun{f}x$ is a non-least element. If also $P_X$ is a
meet-semilattice and $P_Y$ is separable, then
\[
 \supfun{f}(x\sqcap\dom f)=\supfun{f}x .
\]
The second assertion and its proof also
apply to a pair satisfying \eqref{eq-osa-compatible-pair} whenever its
reverse map is monotone.
\end{prop}

\begin{proof}
Intersection with $\top^{P_Y}$ detects non-least elements. Thus the first
assertion follows from
$x\nasymp\supfun{f^\dagger}\top^{P_Y}$ iff
$\top^{P_Y}\nasymp\supfun{f}x$. For the second assertion,
$\supfun{f^\dagger}y\sqsubseteq\dom f$, so
$(x\sqcap\dom f)\nasymp\supfun{f^\dagger}y$ iff
$x\nasymp\supfun{f^\dagger}y$. Apply intersection symmetry and
separability of $P_Y$. For a pair of maps, use $\beta$ in place of
$\supfun{f^\dagger}$ in the same argument.
\end{proof}

\begin{cor}\label{cor-osa-composite-image}
Under the meet and separability hypotheses of Proposition~\ref{prop-osa-support}
for $g:Y\to Z$, if $\im f\sqsupseteq\dom g$ for $f:X\to Y$, then
$\im(g\circ f)=\im g$ whenever both images exist. Under the
corresponding hypotheses for $f^\dagger$, if $\im f\sqsubseteq\dom g$,
then $\dom(g\circ f)=\dom f$ whenever both domains exist.
\end{cor}

\begin{proof}
The image is $\supfun{g}(\im f)=\supfun{g}(\dom g)=\supfun{g}\top^{P_Y}$,
using the support formula twice. Apply this result to
$f^\dagger\circ g^\dagger$ for the domain assertion.
\end{proof}

\begin{prop}\label{prop-osa-pair-joins}
In an intersection-symmetric action, the action maps preserve least
elements when both posets have them. If $P_X,P_Y$ are starrish join-semilattices
and $P_Y$ is separable, then $\supfun{f}$ preserves binary joins.
For a starrish poset (Section~\osaprereqref{sec-filtrator-starrish}),
intersection with a binary join is equivalent to intersection with
at least one of its summands. These assertions
also hold for maps satisfying \eqref{eq-osa-compatible-pair}.
\end{prop}

\begin{proof}
For any $y\in P_Y$, $y\nasymp\supfun{f}\bot^{P_X}$ would imply
$\bot^{P_X}\nasymp\supfun{f^\dagger}y$. Hence the image of the least
element intersects no element, and is therefore least. For $u,v\in P_X$,
intersection of $y$ with $\supfun{f}(u\sqcup v)$ is equivalent to
intersection of $\supfun{f^\dagger}y$ with $u\sqcup v$, hence with
$u$ or $v$. By symmetry and the corresponding starrish law in $P_Y$,
this is equivalent to $y\nasymp(\supfun{f}u\sqcup\supfun{f}v)$.
Separability gives the asserted equality. The proof uses only the
intersection equivalence, so it applies to the indicated pairs as well.
\end{proof}

\begin{defn}\label{def-osa-complete}
\index{completeness!of an interspace}
An interspace $f:X\to Y$ is \emph{complete for the action} if
\[
 \supfun{f}\bigsqcup S=\bigsqcup\rsupfun{\supfun{f}}S
 \qquad(S\subseteq P_X)
\]
whenever both joins exist. If $Z_X\subseteq P_X$ and $Z_Y\subseteq P_Y$
are specified cores of filtrators (Chapter~\osaprereqref{chap-filt}), $f$ is \emph{co-complete relative to
the cores} if $\supfun{f}x\in Z_Y$ for every $x\in Z_X$.
Both predicates apply to a forward map of a pair satisfying
\eqref{eq-osa-compatible-pair} as well.
\end{defn}

For example, these conditions concern joins of elements of $P_X$.
A separate condition concerns joins of arrows. The action
\emph{preserves a join of parallel arrows} $H\subseteq\Hom(X,Y)$ if
$\supfun{\bigsqcup H}x=\bigsqcup_{h\in H}\supfun{h}x$
for every $x\in P_X$, with all indicated joins defined.
In this case evaluation at $\top^{P_X}$ shows that image preserves that
join. If $f:Y\to Z$ acts by a map preserving the resulting joins in
$P_Y$ and the action both preserves the join of $\{f\circ h:h\in H\}$
and is ordered by elements on $\Hom(X,Z)$, then
\[
 f\circ\bigsqcup H=\bigsqcup_{h\in H}(f\circ h).
\]
Indeed, the action maps on both sides agree at every $x\in P_X$ by
the two join preservation hypotheses and the composition law.

\begin{defn}\label{def-osa-element-category}
For an action on a category that preserves identities, the
\emph{category of elements of the action} has objects $(X,x)$
with $x\in P_X$. An arrow from $(X,x)$ to $(Y,y)$ is an
interspace $f:X\to Y$ with $\supfun{f}x\sqsubseteq y$, together
with its indicated source and destination. Composition and
identities use the underlying category.
\end{defn}

Indeed, if $\supfun{g}y\sqsubseteq z$ as well, monotonicity
gives $\supfun{g\circ f}x\sqsubseteq\supfun{g}y\sqsubseteq z$.
The identity sends $x$ to itself.

\section[Restrictions and rectangles]{Restricted identities, restrictions, and rectangles}
\label{sec-osa-restrictions}

\begin{defn}\label{def-osa-restriction}
\index{restricted identity}
An \emph{action with restricted identities} has meet-semilattices $P_X$
and a specified interspace $\id_p:X\to X$ for each $p\in P_X$,
called the \emph{restricted identity} at $p$. They satisfy
\[
\begin{gathered}
 \supfun{\id_p}x=p\sqcap x,\qquad
 \id_q\circ\id_p=\id_{p\sqcap q},\\
 f\circ\id_p\sqsubseteq f,\qquad \id_r\circ f\sqsubseteq f
\end{gathered}
\]
for $p,q,x\in P_X$, $f:X\to Y$, and $r\in P_Y$. In a dagger
action with restricted identities require $\id_p^\dagger=\id_p$.
If the underlying category has identities, the action preserves them,
and $P_X$ has a greatest element, require $\id_{\top^{P_X}}=1_X$.
\end{defn}

The values at $p$ and at $q$ show that the maps
$x\mapsto p\sqcap x$ and $x\mapsto q\sqcap x$ agree only if $p=q$.
In an action ordered by elements the composition law for these maps
and both restriction inequalities already follow from their action.

For $f:X\to Y$, $p\in P_X$ and $q\in P_Y$, its \emph{restriction}
is $f|_p=f\circ\id_p$ and its \emph{rectangular restriction} is
$\id_q\circ f\circ\id_p$. If $X=Y$ and $q=p$, this is its
\emph{square restriction}. These are arrows of the same type as $f$.
Associativity gives $(g\circ f)|_p=g\circ(f|_p)$ for $g:Y\to Z$.
The action and dagger give
\[
 \im(f|_p)=\supfun{f}p,\qquad
 \dom(f|_p)=p\sqcap\dom f
\]
whenever the images or domains in the respective formulas are defined.

\begin{defn}\label{def-osa-rectangle}
\index{product!rectangle}
An ordered action is \emph{correctly bounded} if the posets $P_X$ and
the Hom-sets are bounded, the least arrow acts by the constant least
map, and the greatest arrow $\top_{X,Y}:X\to Y$ takes every non-least
element to $\top^{P_Y}$ and the least element to $\bot^{P_Y}$.
In an action with restricted identities and greatest arrows, the
\emph{binary product} or \emph{rectangle} of $p\in P_X$ and
$q\in P_Y$ is
\[
 p\times q=\id_q\circ\top_{X,Y}\circ\id_p:X\to Y .
\]
The product symbol here denotes an arrow formed from two poset
elements.
\end{defn}

\begin{prop}\label{prop-osa-rectangle-action}
In a correctly bounded action with restricted identities,
\[
 \supfun{p\times q}x=
 \begin{cases}
 q & \text{if }x\nasymp p,\\
 \bot^{P_Y} & \text{otherwise}.
 \end{cases}
\]
Consequently $x\suprel{p\times q}y$ iff $x\nasymp p$ and $y\nasymp q$.
If there is a dagger, $(p\times q)^\dagger=q\times p$.
\end{prop}

\begin{proof}
The three maps in the definition send $x$ successively to $p\sqcap x$,
to $\top^{P_Y}$ or $\bot^{P_Y}$ according as $p\sqcap x$ is non-least
or least, and then to its meet with $q$. For the dagger formula,
the dagger exchanges the greatest arrows in the two Hom-sets and
fixes the restricted identities.
\end{proof}

\begin{prop}\label{prop-osa-rectangle-meet}
Suppose an intersection-symmetric correctly bounded action has
restricted identities, separable posets and order by elements.
Then $f=\id_{\im f}\circ f\circ\id_{\dom f}$, and
\[
 f\sqsubseteq p\times q
 \quad\Longleftrightarrow\quad
 \dom f\sqsubseteq p\ \text{and}\ \im f\sqsubseteq q .
\]
Moreover, $f\sqcap(p\times q)$ exists and equals
$\id_q\circ f\circ\id_p$.
\end{prop}

\begin{proof}
Proposition~\ref{prop-osa-support} makes restriction at $\dom f$
act as $f$ does, and meeting its values with $\im f$ leaves them
unchanged. Order by elements gives the first equality.
If $f\sqsubseteq p\times q$, image and domain monotonicity and
Proposition~\ref{prop-osa-rectangle-action} give the two bounds.
Conversely, those bounds imply $f=\id_q\circ f\circ\id_p$,
using $\id_q\circ\id_{\im f}=\id_{\im f}$ and
$\id_{\dom f}\circ\id_p=\id_{\dom f}$. This implies
$f\sqsubseteq p\times q$. For the meet formula the three-factor
composite lies below both $f$ and $p\times q$. Any common lower
bound $h$ satisfies $h=\id_q\circ h\circ\id_p$ by the same bounds
and hence lies below the composite.
\end{proof}

The proof of the meet formula also applies to actions in which
$f=\id_{\im f}\circ f\circ\id_{\dom f}$ has been established by a
different argument. In a correctly bounded, intersection-symmetric
action ordered by elements, Proposition~\ref{prop-osa-support} also
shows that $\im f$ is least iff $f$ is the least arrow. Thus in the
setting of Proposition~\ref{prop-osa-rectangle-meet},
$f\nasymp(p\times q)$ iff $p\suprel{f}q$. In fact, the image of the
meet is $q\sqcap\supfun{f}p$ by the restriction formula.

\begin{prop}\label{prop-osa-atomic-rectangles}
In the setting of Proposition~\ref{prop-osa-rectangle-meet},
if $P_X$ and $P_Y$ are atomic, the rectangles of two atoms
are exactly the atoms of $\Hom(X,Y)$. If $P_X$ and $P_Y$
are atomistic, every $f:X\to Y$ is the join of those atom
rectangles below it.
\end{prop}

\begin{proof}
Let $p\in P_X$ and $q\in P_Y$ be atoms. If $g$ is non-least and
$g\sqsubseteq p\times q$, the domain and image bounds force
$\dom g=p$ and $\im g=q$. Both are non-least, since a least
image or domain forces the arrow to be least. By the domain
test in Proposition~\ref{prop-osa-support}, $\supfun{g}x$ is
non-least iff $x\nasymp p$, and then its bound $q$ forces
$\supfun{g}x=q$. Thus $g=p\times q$ by order by elements.
Conversely, any non-least $f$ has a nonempty intersection
relation, so by Proposition~\ref{prop-osa-atom-recovery} there
are atoms $p,q$ with $p\suprel{f}q$. The meet $f\sqcap(p\times q)$
is non-least by the intersection test for rectangles and
hence equals $p\times q$. This proves the characterization of
atoms.

When both posets are atomistic, each atom pair related by $f$
gives a rectangle below $f$. Any upper bound of these rectangles
relates all those pairs, so it lies above $f$ by the order test in
Proposition~\ref{prop-osa-atom-recovery}. This gives the join.
\end{proof}

The two case formulas for $p\times q$ and $q\times p$ also define
a compatible pair on any two posets with least elements, without
requiring greatest elements or restricted identities. In the
compatibility condition for this pair, both sides hold exactly
when $x\nasymp p$ and $y\nasymp q$. In a correctly bounded
intersection-symmetric action with restricted identities, these
maps are the action of the rectangle and its dagger.

In the same setting the action on rectangles satisfies
\[
 g\circ(p\times q)\circ h=\supfun{h^\dagger}p\times\supfun{g}q
\]
for $h:W\to X$ and $g:Y\to Z$. To see this, first compose on the
right. By intersection symmetry, $\supfun{h}w\nasymp p$ iff
$w\nasymp\supfun{h^\dagger}p$. Now compose on the left and use
preservation of the least element by $\supfun{g}$. The resulting map
is the action of the rectangle on the right, so order by elements
gives the identity.

\begin{defn}\label{def-osa-restricted-category}
Let an action with restricted identities have an underlying category.
Its \emph{category of restricted interspaces} has objects $(X,p)$
with $p\in P_X$. An arrow from $(X,p)$ to $(Y,q)$ is an interspace
$f:X\to Y$ satisfying $f=\id_q\circ f\circ\id_p$, together
with its specified source and destination pairs. Composition uses
the underlying composition, and the identity of $(X,p)$ is $\id_p$.
\end{defn}

This is a category because $\id_p$ is idempotent. The defining equality
is equivalent to $\id_q\circ f=f=f\circ\id_p$, so it is preserved by
composition. The dagger, when available, exchanges the source and
destination pairs. Under the hypotheses of
Proposition~\ref{prop-osa-rectangle-meet}, the condition on an arrow
is equivalent to $\dom f\sqsubseteq p$ and $\im f\sqsubseteq q$.
Indeed, that pair of bounds implies the defining equality, and the
equality implies $f\sqsubseteq p\times q$. For a self-adjoint
$k:Y\to Y$ in that setting, $k\sqsubseteq 1_Y$ implies
$k\sqsubseteq\id_{\im k}$, by restricting both sides at $\im k$.
In particular, for $k=f\circ f^\dagger$ one has
$\im k\sqsubseteq\im f$. Restricting at $\im f$ gives the equivalent
form $f\circ f^\dagger\sqsubseteq\id_{\im f}$ of monovaluedness.
The dagger gives the corresponding form of injectivity using $\dom f$.
Here monovaluedness and injectivity use the identity inequalities of
Section~\osaprereqref{sec-ordered-morphism-classes}.

\begin{prop}\label{prop-osa-monovalued-atoms}
Let an intersection-symmetric action preserve categorical identities
and be ordered by elements. If $P_X$ and $P_Y$ are atomistic
meet-semilattices with least elements, the following are equivalent
for $f:X\to Y$.
\begin{enumerate}
\item $f$ is monovalued.
\item $\supfun{f}p$ is least or an atom for every atom $p$ of $P_X$.
\item $\supfun{f^\dagger}$ preserves binary meets.
\end{enumerate}
\end{prop}

\begin{proof}
On atoms $q,r$ of $P_Y$, the relation of $f\circ f^\dagger$ relates
$q$ to $r$ precisely when some atom $p$ of $P_X$ is related by $f$
to both of them, by Proposition~\ref{prop-osa-rel-comp}. The identity
relates only equal atoms, so the order test in
Proposition~\ref{prop-osa-atom-recovery} gives the equivalence of the
first two conditions. For the second to imply the third, test
$\supfun{f^\dagger}(u\sqcap v)$ on a source atom $p$.
The test is whether $\supfun{f}p$ intersects $u\sqcap v$, which
for a least element or an atom is equivalent to intersecting both
$u$ and $v$. That is exactly the test for
$\supfun{f^\dagger}u\sqcap\supfun{f^\dagger}v$. For the converse,
distinct atoms $q,r$ in $P_Y$ have least meet. If their inverse
images have least meet, no source atom can be related to both.
Use preservation of the least element by the reverse map and
apply the first equivalence.
\end{proof}

\section{Continuity and anticontinuity}\label{sec-osa-continuity}

\begin{defn}\label{def-osa-continuity}
\index{continuity!of interspaces}\index{anticontinuity}
Let $\mu:X\to X$ and $\nu:Y\to Y$ be endointerspaces. The set
$\continuous(\mu,\nu)$ consists of the arrows $f:X\to Y$ satisfying
$f\circ\mu\sqsubseteq\nu\circ f$. In a dagger action, the sets
$\continuous'(\mu,\nu)$ and $\continuous''(\mu,\nu)$ consist,
respectively, of the arrows satisfying
\[
 \mu\sqsubseteq f^\dagger\circ\nu\circ f,
 \qquad
 f\circ\mu\circ f^\dagger\sqsubseteq\nu .
\]
These are the three forms of \emph{continuity}. Their
\emph{anticontinuity} counterparts
$\continuous_\ast(\mu,\nu)$, $\continuous_\ast'(\mu,\nu)$,
and $\continuous_\ast''(\mu,\nu)$ use the reverse inequality
in each corresponding formula.
\end{defn}

Each continuity formula gives a semicategory of endointerspaces
with the indicated arrows between them. When the underlying
semicategory has identities it gives a category. For the first
formula, composing $f\in\continuous(\mu,\nu)$ and
$g\in\continuous(\nu,\pi)$ gives
$g\circ f\circ\mu\sqsubseteq g\circ\nu\circ f\sqsubseteq
\pi\circ g\circ f$. The corresponding chains for the other two are
\[
\begin{gathered}
 \mu\sqsubseteq f^\dagger\circ\nu\circ f
 \sqsubseteq (g\circ f)^\dagger\circ\pi\circ(g\circ f),\\
 (g\circ f)\circ\mu\circ(g\circ f)^\dagger
 \sqsubseteq g\circ\nu\circ g^\dagger\sqsubseteq\pi.
\end{gathered}
\]
The identity arrows satisfy each formula. Reversing the order gives
the same conclusions for anticontinuity. The dagger carries each
anticontinuity condition to a continuity condition for
$f^\dagger$ from $\nu^\dagger$ to $\mu^\dagger$, leaving the
unprimed form in place and exchanging the two primed forms.

\begin{defn}\label{def-osa-cont-category}
For an ordered category $\mathcal C$, write
$\operatorname{cont}(\mathcal C)$ for the category with objects
the endomorphisms of $\mathcal C$ and arrows $(f,\mu,\nu)$ with
$f\in\continuous(\mu,\nu)$. Composition is
$(g,\nu,\pi)\circ(f,\mu,\nu)=(g\circ f,\mu,\pi)$ and the
identity of $\mu:X\to X$ is $(1_X,\mu,\mu)$.
If only arrows $f$ in a specified subcategory $\mathcal D$ with
all objects of $\mathcal C$ are allowed, write
$\operatorname{cont}(\mathcal C;\mathcal D)$. It is a subcategory
with all the same endomorphisms as objects.
\end{defn}

\begin{thm}\label{cont-eq}
In an ordered dagger action with categorical identities, the three
continuity conditions satisfy
$\continuous'\Rightarrow\continuous\Rightarrow\continuous''$ on
monovalued interspaces and the reverse implications on entirely defined
interspaces. Thus they coincide on interspaces with both properties.
On such interspaces,
$f\in\continuous(\mu,\nu)$ iff
$f\in\continuous(\mu^\dagger,\nu^\dagger)$.
\end{thm}

\begin{proof}
For a monovalued $f$, compose the primed inequality on the left by $f$
and use $f\circ f^\dagger\sqsubseteq1_Y$ to get the unprimed
inequality. Compose the latter on the right by $f^\dagger$ and use
the same bound to get the double-primed inequality. For an entirely
defined $f$, use $1_X\sqsubseteq f^\dagger\circ f$: compose the
double-primed inequality on the right by $f$, or the unprimed
inequality on the left by $f^\dagger$, respectively.
Applying the dagger to the primed inequality shows that it holds
exactly when the primed inequality for $\mu^\dagger,\nu^\dagger$
holds. The equivalence of all three forms gives the last assertion.
\end{proof}

\begin{thm}\label{rect-cont}
For an action with restricted identities, let $p\in P_X$ and
$\mu_p=\id_p\circ\mu\circ\id_p$. Then
\begin{enumerate}
\item\label{contrestr-C}
$f\in\continuous(\mu,\nu)\Rightarrow f|_p\in\continuous(\mu_p,\nu)$.
\item\label{contrestr-Ci}
In a dagger action,
$f\in\continuous'(\mu,\nu)\Rightarrow f|_p\in\continuous'(\mu_p,\nu)$.
\item\label{contrestr-Cii}
In a dagger action,
$f\in\continuous''(\mu,\nu)\Rightarrow f|_p\in\continuous''(\mu_p,\nu)$.
\end{enumerate}
\end{thm}

\begin{proof}
For the unprimed case use idempotence of $\id_p$ and the restriction
inequality to obtain
$f|_p\circ\mu_p=f\circ\id_p\circ\mu\circ\id_p
\sqsubseteq f\circ\mu\circ\id_p\sqsubseteq\nu\circ f|_p$.
For the primed case compose the original inequality on both sides by
$\id_p$. For the double-primed case, the new left-hand side is
$f\circ\id_p\circ\mu\circ\id_p\circ f^\dagger$, which lies
below $f\circ\mu\circ f^\dagger$.
\end{proof}

\section[Closed elements and convergence]{Neighborhoods, closed elements, and convergence}
\label{sec-osa-convergence}

\begin{defn}\label{def-osa-convergence}
\index{convergence!in an ordered action}
Given $\mu:X\to Y$ and $a\in P_X$, an element $b\in P_Y$ is a
\emph{neighborhood} of $a$ regarding $\mu$ if
$\supfun{\mu}a\sqsubseteq b$. The \emph{interior} of $b$ regarding
$\mu$ is $\bigsqcup\setcond{a\in P_X}{\supfun{\mu}a\sqsubseteq b}$
when that join exists. If $\supfun{\mu}$ preserves this join, the
interior is the greatest element having $b$ as a neighborhood.

An element $v\in P_Y$ \emph{converges} to $a$ regarding $\mu$,
written $v\overset{\mu}{\longrightarrow}a$, if
$v\sqsubseteq\supfun{\mu}a$. An interspace $f:W\to Y$
\emph{converges} to $a$ regarding $\mu$ if its image is defined
and $\im f\overset{\mu}{\longrightarrow}a$. It \emph{converges
on $p$} for $p\in P_W$ if
$\supfun{f}p\overset{\mu}{\longrightarrow}a$. With restricted
identities and images this is equivalent to convergence of $f|_p$.
\end{defn}

\begin{defn}\label{def-osa-closed}
\index{closed element}\index{open element}
For an endointerspace $\mu:X\to X$, an element $x\in P_X$ is
\emph{closed regarding $\mu$} if $\supfun{\mu}x\sqsubseteq x$.
In a dagger action, $x$ is \emph{open regarding $\mu$} if it is
closed regarding $\mu^\dagger$. Given another endointerspace
$\nu:Y\to Y$, an interspace $f:X\to Y$ is \emph{open regarding
$\mu$ and $\nu$} if $\supfun{f}$ sends elements open regarding $\mu$
to elements open regarding $\nu$. The endointerspace $\mu$ is
\emph{closed} if every element is closed regarding it.
When the action preserves the
identity, $\mu\sqsubseteq1_X$ implies that every element is closed
regarding $\mu$. In an action ordered by elements the converse holds.
\end{defn}

Any existing meet of closed elements is closed, since monotonicity
places the action on the meet below each member of the family.
A join of closed elements is closed whenever the action map preserves
that join. In particular, a complete endointerspace on a complete
lattice has both closure properties. I call an idempotent endointerspace
that is complete for the action a \emph{Kuratowski space}. For a
functional action on a powerset, requiring also $1_X\sqsubseteq\mu$
makes it a topological closure operator when all joins are preserved.

\begin{prop}\label{prop-osa-cont-conv}
Let $f:W\to X$, $g:X\to Y$, $\mu:X\to X$, and $\nu:Y\to Y$.
If $g\in\continuous(\mu,\nu)$, then
$f\overset{\mu}{\longrightarrow}a$ implies
$g\circ f\overset{\nu}{\longrightarrow}\supfun{g}a$ whenever
$\im f$ is defined. The same implication holds for convergence on
$p\in P_W$. Continuity implies the pointwise inequalities
\[
 \supfun{g}\supfun{\mu}a\sqsubseteq\supfun{\nu}\supfun{g}a
 \qquad(a\in P_X).
\]
In an action ordered by elements these inequalities imply continuity.
\end{prop}

\begin{proof}
Apply the action to the continuity inequality. For convergence of
arrows, use
$\im(g\circ f)=\supfun{g}(\im f)\sqsubseteq
\supfun{g}\supfun{\mu}a\sqsubseteq\supfun{\nu}\supfun{g}a$.
For convergence on $p$ use $\supfun{f}p$ in place of $\im f$.
The converse follows by order by elements.
\end{proof}

The composition result also has a local form. Let
$b=\supfun{\mu}a$ and $a\sqsubseteq b$. In an action with restricted
identities, write $\mu_b=\id_b\circ\mu\circ\id_b$. If
$g|_b\in\continuous(\mu_b,\nu)$ and
$f\overset{\mu}{\longrightarrow}a$, then again
$g\circ f\overset{\nu}{\longrightarrow}\supfun{g}a$.
Indeed, $\im f\sqsubseteq b$ makes
$\im(g\circ f)=\supfun{g|_b}(\im f)$. Since
$\supfun{\mu_b}a=b$, local continuity bounds
$\supfun{g|_b}b$ by $\supfun{\nu}\supfun{g|_b}a$,
which equals $\supfun{\nu}\supfun{g}a$.

There is also a local form in an action with restricted identities.
For $\mu:X\to X$, $g:X\to Y$, and $a\in P_X$, the inequality
$\supfun{g\circ(\mu|_a)}a\sqsubseteq\supfun{\nu\circ g}a$
is equivalent to convergence of $g$ to $\supfun{g}a$ regarding
$\nu:Y\to Y$ on $\supfun{\mu}a$. This uses
$\supfun{\mu|_a}a=\supfun{\mu}a$.
When $a$ is an atom, the poset $P_X$ has a least element,
$\supfun{g}$ and $\supfun{\mu}$ preserve it, and the action is
ordered by elements, the local inequality is equivalent to
$g\in\continuous(\mu|_a,\nu)$. Indeed, for $x\in P_X$, the meet
$a\sqcap x$ is either least or $a$. In the former case the
left-hand action value is least; in the latter use the local
inequality and $a\sqsubseteq x$.

For a family $V\subseteq P_Y$ with an existing join, $\bigsqcup V$
converges to $a$ regarding $\mu:X\to Y$ iff each member of $V$ does.
The same equivalence holds for a join of parallel interspaces whose
images preserve the join. For convergence of $f$ on a join of
elements of $P_W$, the equivalence holds if $\supfun{f}$ preserves
that join. In all cases, use the least-upper-bound property with
upper bound $\supfun{\mu}a$.

\section{Separability, limits, and compactness}
\label{sec-osa-separation}

\begin{defn}\label{def-osa-separation}
\index{separation!of interspaces}
Choose a subset of pairwise nonintersecting non-least elements of each
poset of interest and call the chosen elements \emph{points}. For
$\mu:X\to X$, call $\mu$ \emph{$T_1$-separable} if for points $p,q$
in $P_X$, $p\suprel{\mu}q$ implies $p=q$. In a dagger action,
$\mu$ is \emph{$T_0$-separable} if the meet
$\mu\sqcap\mu^\dagger$ exists and is $T_1$-separable. An arrow
$\mu:X\to Y$ is \emph{$T_2$-separable} if $\mu^\dagger\circ\mu$
is $T_1$-separable. An endointerspace $\mu:X\to X$ is
\emph{$T_4$-separable} if
\[
 (\mu\circ\mu^\dagger)\circ(\mu\circ\mu^\dagger)
 \sqsubseteq \mu\circ\mu^\dagger.
\]

Given also a subset $Z_Y\subseteq P_Y$ of test elements,
$\mu:X\to Y$ is \emph{regular} relative to the points and $Z_Y$ if
the following implication holds for every point $p$ of $P_X$
and every $z\in Z_Y$.
\[
 z\nasymp\supfun{\mu\circ\mu^\dagger\circ\mu}p
 \quad\Longrightarrow\quad z\nasymp\supfun{\mu}p.
\]
An endointerspace is \emph{$T_3$-separable} if it is $T_2$-separable
and regular for the specified test elements.
For a symmetric idempotent endointerspace, all three conditions
$T_0$, $T_1$, and $T_2$ coincide because the meet and composite
in their definitions both equal the endointerspace itself.
\end{defn}

In an intersection-symmetric action, the composition formula in
Proposition~\ref{prop-osa-rel-comp} makes $T_2$-separability exactly
the condition that the values $\supfun{\mu}p$ and $\supfun{\mu}q$
do not intersect for distinct points $p,q$.

\begin{defn}\label{def-osa-limit}
A point $p$ of $P_X$ is a \emph{limit} of $v\in P_Y$ regarding
$\mu:X\to Y$ if $v\overset{\mu}{\longrightarrow}p$. The limits
of an interspace $f:W\to Y$ are the limits of its image,
when that image exists. Its limits on $a\in P_W$ are the limits
of $\supfun{f}a$. When there is exactly one limit use the
notations $\lim^\mu f$ and $\lim_a^\mu f$ for the last two cases.
\end{defn}

In an intersection-symmetric action, if $\mu$ is $T_2$-separable
and $v$ is non-least, there is at most one limit of $v$. If there
were distinct limits $p,q$, then $v$ would witness intersection of
$\supfun{\mu}p$ and $\supfun{\mu}q$.

\begin{defn}\label{def-osa-xlim}
\index{limit!generalized}
Suppose the underlying semicategory is a category. Fix a group $G$,
a homomorphism $\rho:G\to\operatorname{Aut}(X)$ to the group of
invertible arrows $X\to X$, and an endointerspace $\nu:Y\to Y$.
The \emph{generalized limit} of $f:X\to Y$ for these choices is
the subset of $\Hom(X,Y)$
\[
 \xlim f=\setcond{\nu\circ f\circ\rho(r)}{r\in G}.
\]
If $\mu:X\to X$ is given and the action has restricted identities,
define the generalized limit at $a\in P_X$ by
$\xlim_a f=\xlim(f|_{\supfun{\mu}a})$.
\end{defn}

\begin{defn}\label{def-osa-compact}
\index{compactness!of an interspace}
Let $f:X\to Y$, where the destination action poset is a filtrator
$(P_Y,Z_Y)$. Assume its core part $\Cor y$ is defined for each
$y\in P_Y$ and both $Z_Y$ and $P_Y$ have least elements.
Then $f$ is \emph{directly compact} if for each $x\in P_X$,
\[
 \supfun{f}x\ne\bot^{P_Y}
 \quad\Longrightarrow\quad
 \Cor\supfun{f}x\ne\bot^{Z_Y}.
\]
In a dagger action, $f$ is \emph{reversely compact} if
$f^\dagger$ is directly compact using the source filtrator.
It is \emph{compact} if both conditions hold.
Here $\Cor$ is the core part of Section~\osaprereqref{sec-filtrator-core-part}.
\end{defn}

If the domain of $f$ is defined in an intersection-symmetric
action with atomic source poset having a least element, direct
compactness can be tested on atoms under $\dom f$. In fact, every
$x$ with non-least image intersects $\dom f$, so it has an atom
$a\sqsubseteq x,\dom f$. The image of $a$ is non-least, and
monotonicity of $\supfun{f}$ and $\Cor$ extends the condition at $a$
to $x$. The reverse implication follows from testing each such atom.

\section{Connectedness and endomorphism series}
\label{sec-osa-connected}

\begin{defn}\label{def-osa-connected}
\index{connectedness!in an ordered action}
Let $P_X$ have a least element and let $\mu:X\to X$. An element
$a\in P_X$ is \emph{connected regarding $\mu$} in the sense of its
intersection relation if, for non-least $x,y\in P_X$,
\[
 x\sqcup y=a\quad\Longrightarrow\quad x\suprel{\mu}y .
\]
The condition tests pairs whose join exists.
\end{defn}

\begin{defn}\label{def-osa-series}
\index{endomorphism series}
For $\mu:X\to X$, write $\mu^1=\mu$ and
$\mu^{n+1}=\mu\circ\mu^n$ for positive integers $n$. The
\emph{endomorphism series} are
\[
 S_1(\mu)=\bigsqcup_{n\ge1}\mu^n,
 \qquad
 S(\mu)=\bigsqcup_{n\ge0}\mu^n,
\]
when the indicated joins exist in $\Hom(X,X)$ and, for the second
formula, $\mu^0=1_X$ exists. Whenever composition distributes over the
joins in the second formula, $S(\mu)$ is reflexive and idempotent.
Indeed, the composite of the two joins is the join of the powers
$\mu^{m+n}$ with $m,n\ge0$, which gives the same join.

Suppose now that a subset $Z\subseteq\Hom(X,X)$ is specified,
and that the series are defined on $Z$. Write
$\up_Z\mu=\setcond{e\in Z}{\mu\sqsubseteq e}$.
Define the \emph{connectivity series relative to $Z$} by
\[
 S_1^\ast(\mu)=\bigsqcap_{e\in\up_Z\mu} S_1(e),\qquad
 S^\ast(\mu)=\bigsqcap_{e\in\up_Z\mu} S(e),
\]
where the meets are in $\Hom(X,X)$. Given rectangles,
$a\in P_X$ is \emph{connected regarding $\mu$ relative to $Z$} in
this sense if the following meet and series exist and
\[
 a\times a\sqsubseteq S_1^\ast(\mu\sqcap(a\times a)).
\]
\end{defn}

If $Z$ is closed under $S$, and $S$ is extensive and idempotent on
$Z$, then $S^\ast$ is idempotent wherever both applications are
defined. In fact $S^\ast(\mu)\sqsupseteq\mu$. For each $e\in\up_Z\mu$,
$S(e)$ belongs to $\up_Z S^\ast(\mu)$. Hence
$S^\ast(S^\ast(\mu))\sqsubseteq S(S(e))=S(e)$ for each $e$.
Taking the meet gives the reverse of the inequality from extensivity.

\section{Rectangles and filtrators}
\label{sec-osa-rectangles-topology}

\begin{defn}\label{def-osa-rectangle-complete}
In an action with rectangles, choose points in the action posets
and use their atoms. The set of atoms can be empty. An interspace $f:X\to Y$ is
\emph{complete relative to these rectangles} if it is the join
of a family of rectangles $p\times b$ with $p$ a chosen point of
$P_X$ and $b$ an atom of $P_Y$. It is \emph{co-complete relative
to these rectangles} if it is a join of rectangles $a\times q$
with $a$ an atom of $P_X$ and $q$ a chosen point of $P_Y$.
In a dagger action the dagger exchanges these two predicates,
by Proposition~\ref{prop-osa-rectangle-action}.
\end{defn}

\begin{defn}\label{def-osa-completion}
\index{completion!of an interspace}
Given a specified class of complete arrows in each Hom-set
(for example, Definition~\ref{def-osa-complete} or
\ref{def-osa-rectangle-complete}), the \emph{completion} $\Compl f$ is
the greatest complete arrow below $f$ in the Hom-set of $f$, when
it exists. With a specified class of co-complete arrows, the
\emph{co-completion} $\CoCompl f$ is the greatest co-complete arrow
below $f$, when it exists. If the co-complete arrows are the
daggers of the complete arrows, then
$\CoCompl f=(\Compl(f^\dagger))^\dagger$. If the complete arrows below $f$ have a join which is
complete, that join is the completion. In particular, for
completeness relative to rectangles, the join of all the specified
generator rectangles below $f$ is the completion when it exists.
\end{defn}

\begin{defn}\label{def-osa-thickness}
\index{total boundedness}
Suppose $P_X$ has a greatest element and has finite joins.
Choose a subset $K\subseteq P_X$ of allowed covering elements,
with $K=P_X$ if no restriction is specified.
For an endointerspace $e:X\to X$ with rectangles, write
$\thick_\alpha(e)$, or call it \emph{$\alpha$-thick}, if some
finite $A\subseteq K$ has $\bigsqcup A=\top^{P_X}$ and
$a\times a\sqsubseteq e$ for all $a\in A$. Given a set of points
in $P_X$, write $\thick_\beta(e)$, or call it \emph{$\beta$-thick},
if for some finite set $B$ of those points,
$\supfun{e}(\bigsqcup B)=\top^{P_X}$.
Relative to $Z\subseteq\Hom(X,X)$, an endointerspace $\mu$
is \emph{$\alpha$-totally bounded} (respectively
\emph{$\beta$-totally bounded}) if every $e\in\up_Z\mu$ is
$\alpha$-thick (respectively $\beta$-thick).
\end{defn}

In a correctly bounded action with restricted identities, if every
non-least element contains a chosen point below it, then
$\alpha$-thickness implies $\beta$-thickness. For a finite family $A$ with join $\top^{P_X}$, choose a point $p_a\sqsubseteq a$ below each
non-least $a$. Since $a\times a\sqsubseteq e$,
Proposition~\ref{prop-osa-rectangle-action} gives
$a\sqsubseteq\supfun{e}p_a$. Monotonicity gives the same bound for
$\supfun{e}(\bigsqcup_a p_a)$, so that value is greatest.

A \emph{lower set} in an action poset is a subset $C\subseteq P_X$
such that $q\in C$ whenever $q\sqsubseteq p\in C$.
If $P_X$ has a least element, I call such a family \emph{totally
bounded in the Cauchy sense} if every non-least $p\in P_X$ has a
non-least $q\in C$ with $q\sqsubseteq p$. This condition concerns
a lower set, whereas Definition~\ref{def-osa-thickness} concerns
upper approximants of an interspace.

\begin{defn}\label{def-osa-cauchy}
\index{Cauchy completeness}
In an action with rectangles, an element $p\in P_X$ is
\emph{Cauchy regarding $\mu:X\to X$} if $p\times p\sqsubseteq\mu$.
With chosen points, $\mu$ is \emph{Cauchy-complete} if every
Cauchy element converges to a point regarding $\mu$.
For specified subsets $C_X\subseteq P_X$ and $C_Y\subseteq P_Y$
of Cauchy elements of interest, an arrow $f:X\to Y$ is
\emph{Cauchy continuous} from $C_X$ to $C_Y$ if
$\supfun{f}p\in C_Y$ for every $p\in C_X$.
\end{defn}

For a subset $C\subseteq P_X$, the join
$\bigsqcup_{p\in C}(p\times p)$, when it exists, is the least
endointerspace regarding which every element of $C$ is Cauchy.
This follows from the defining inequality for Cauchy elements
and the least-upper-bound property.

\section[Projection and injection products]{Products specified by projections or injections}
\label{sec-osa-projection-products}

These products use a chosen family of projections or injections
in an ordered dagger semicategory. Let $I$ be a small index set
and let $A=(A_i)_{i\in I}$ and $B=(B_i)_{i\in I}$ be families of
objects. Suppose chosen objects $U_A,U_B$ have projection arrows
$\pi_i^A:U_A\to A_i$ and $\pi_i^B:U_B\to B_i$, respectively.
Here the product objects are specified independently of the action
posets. For a family $F=(F_i:A_i\to B_i)_{i\in I}$ define
\begin{equation}\label{def-osa-projective-product}
 \bigodot^{\mathrm{proj}}_{\sqcap}F=
 \bigsqcap_{i\in I}\bigl((\pi_i^B)^\dagger\circ F_i\circ\pi_i^A\bigr)
 \ \in\Hom(U_A,U_B)
\end{equation}
when the meet exists. The subscript $\sqcup$ specifies the
join of the same family. For a common source $W$ and
$F_i:W\to B_i$, use
$\prod^{\mathrm{proj}}_{\sqcap}F=
\bigsqcap_{i\in I}((\pi_i^B)^\dagger\circ F_i)$ in $\Hom(W,U_B)$.
For a common destination $W$ and $F_i:A_i\to W$, use
$\coprod^{\mathrm{proj}}_{\sqcup}F=
\bigsqcup_{i\in I}(F_i\circ\pi_i^A)$ in $\Hom(U_A,W)$.
The other subscript in each case switches meet and join.
If instead there are chosen objects with injection arrows from
the family members, the superscript $\mathrm{inj}$ uses the
daggers of those injections in place of the indicated projections.

In a category, consider endomorphisms $F_i:A_i\to A_i$ and
projections that are entirely defined and monovalued. The meet
product of \eqref{def-osa-projective-product} is the greatest
endomorphism of the chosen product object for which each
projection is continuous. For any endomorphism $M$ of that
object, by Theorem~\ref{cont-eq}, continuity of $\pi_i^A$ from
$M$ to $F_i$ is equivalent to
$M\sqsubseteq(\pi_i^A)^\dagger\circ F_i\circ\pi_i^A$.
The assertion is the defining property of the meet.
